\begin{abstract}
Sparse surrogate models describe neural activations using a small number
of active latent coordinates. We study how these coordinates can be
combined into reproducible queries about images and spatial sites.
Following a relation-induced construction for language-model features,
we derive a lattice-theoretic formal context from nonnegative visual
codes. Common descriptions are coordinatewise minima, and retrieval is
coordinatewise threshold satisfaction. We retain the Galois-ideal and
adjunction proofs and identify the construction with a pattern structure
and, on finite data, an ordinally scaled classical context. The
vision-specific issue is spatial quantification: max pooling allows
different coordinates to be witnessed at different sites. We characterize
when one vector can represent same-site evidence and give a downset-valued
construction when it cannot. Three digit CNN/autoencoder runs evaluate
held-out retrieval and reconstruction. A frozen ResNet34 with a sparse
autoencoder supplies measured animal and vehicle feature galleries on a
fixed 1,500-image CIFAR-10 subset. Explicit query examples distinguish
pooled matches from same-site witnesses. Exhaustive finite checks verify
the algebra independently. Released code, weights, images, and activations
make the examples auditable; the results do not establish causal or
class-exclusive interpretations of the learned coordinates.
\end{abstract}

\noindent\textbf{Keywords:} formal concept analysis; pattern structures;
computer vision; sparse autoencoders; spatial feature analysis.

\section{Introduction}\label{sec:introduction}
A vision model can recognize an image accurately while leaving unclear
which internal features support its decision. Feature visualization
investigates patterns that activate individual units, whereas sparse
surrogates learn a larger dictionary in which an activation is represented
by relatively few nonzero coefficients
\cite{olah2017featureVis,gao2024,lim2025patchsae}.
Neither construction makes a coordinate a semantic concept by definition.
A useful next question is relational: which images share a collection of
features, at which activation levels, and in which spatial locations?

Formal concept analysis (FCA) organizes objects and their shared
attributes into a concept lattice~\cite{ganter1999formal}. In the present
setting, the objects are images or spatial sites, and an attribute query
specifies minimum values for selected sparse coordinates. For example, a
query requiring two features at levels $2$ and $1$ retrieves precisely the
objects whose codes dominate $(2,1)$. The common description of those
objects is their coordinatewise minimum. Closing the query replaces it
by this common description without changing its retrieved objects.
This is a mathematical account of shared evidence, not a guarantee that
the evidence has the interpretation suggested by a feature visualization.

The paper adapts the summary--retrieval framework of the companion
language-model manuscript~\cite{tsinadzeCompanion}. We preserve its
sequence of arguments: construct a Galois ideal, derive the antitone
Galois connection, identify the context derivations, and relate the
resulting concepts to classical FCA. These arguments are applications of
established theory~\cite{FCA_Latt_Theor,ganter2001pattern}; we make no
novelty claim for the existence of the adjunction or concept lattice.

Images introduce a distinction that cannot be handled by changing the
word ``token'' to ``patch.'' If one patch exhibits the first feature and
another exhibits the second, their coordinatewise maximum satisfies the
two-feature query even though neither patch satisfies it. The difference
is between ``for each coordinate, some site'' and ``some site, for every
coordinate.'' We formalize both readings, characterize the special case
in which a single vector represents them equally, and use a lattice of
downsets to retain same-site evidence in general. Site-level evidence
itself need not be spatially local in the input: attention and convolution
can give a site a receptive field larger than its displayed cell.

Our empirical contribution combines a controlled digit pilot with a
natural-image feature study. The digit pilot uses a convolutional
classifier, a nonnegative Top-$k$ sparse autoencoder, separate training,
calibration, and test images, and saved query vectors and scores. A frozen
pretrained ResNet34 supplies animal and vehicle examples. Public vision
libraries supply a route to further models; their published benchmark
numbers are not results of this paper. Our studies test whether the
proposed operations can be computed and interpreted consistently. They
do not test foundation-model superiority or establish a new state of
the art.

The exposition first specifies sparse surrogates and spatial summaries,
then develops the context and its spatial refinement. We subsequently
report computational checks, the digit benchmark, and natural-image
examples, and conclude with
the scope of the evidence and requirements for a larger evaluation.
